\chapter{Resultados y conclusiones}
\label{cap:conclusiones}

En este capítulo se juntan los resultados de Trabajo Terminal II, se comparan con los
objetivos y requerimientos planteados en TT1 y se presentan las conclusiones del
trabajo. El detalle de cada medición está en el Capítulo~\ref{cap:desarrollo}.

% ==============================================================================
\section{Resumen de resultados}
% ==============================================================================

La Tabla~\ref{tab:resumen_resultados} resume el resultado principal de cada módulo y
las condiciones en las que se obtuvo.

\begin{longtable}{|p{2.6cm}|p{6.4cm}|p{5.5cm}|}
\caption{Resultados principales por módulo}
\label{tab:resumen_resultados} \\
\hline
\textbf{Módulo} & \textbf{Resultado} & \textbf{Condición de la prueba} \\
\hline
\endfirsthead
\hline
\textbf{Módulo} & \textbf{Resultado} & \textbf{Condición de la prueba} \\
\hline
\endhead
Prototipo & \TODO{resultado general de los recorridos con el prototipo} &
\TODO{lugares, participantes y número de recorridos} \\
\hline
Visión & Detecta personas, muebles, puertas con su estado, escaleras y peligros de piso,
en unos 160~ms por cuadro. & 19 videos reales, CPU de computadora portátil;
\TODO{precisión en imágenes etiquetadas del prototipo}. \\
\hline
SLAM & ATE de 0.21~m en recorridos cortos; con odometría sesgada, el error queda entre
la mitad y tres cuartas partes del de la odometría sola. & Banco de simulación, seis
semillas. \\
\hline
Ray Casting & Actualización y consulta del mapa en menos de 1~ms por ciclo. &
Simulación y videos reales. \\
\hline
Mapeo & Mapa por bloques con persistencia; 19 mapas de lugares reales de entre 24 y
60~KB. & Videos reales. \\
\hline
Navegación & 97.8\,\% de llegadas, 2.0\,\% de choques y ninguna caída en 880
episodios. & Simulador de entrenamiento con 19 mapas reales. \\
\hline
Navegación (sistema completo) & 20 de 20 sesiones con llegada; en 17 el usuario terminó
a menos de 1~m del objeto. & Lazo cerrado con todos los módulos y un usuario
virtual. \\
\hline
Alertas & Avisos críticos y de precaución con pitido estéreo; evaluación en 0.007~ms
por ciclo. & Escenario simulado con una persona cruzando. \\
\hline
PLN & 51.2\,\% de intenciones correctas en 86 frases coloquiales que el PLN no conocía,
en 4.8~ms por frase. & Frases redactadas antes de la evaluación y fuera del conjunto de
referencia. \\
\hline
Voz completa & 95.0\,\% de comandos correctos en limpio con voces naturales, y de 62 a
74\,\% con ruido de 5 a 10~dB. & Voz sintetizada, Vosk y PLN;
\TODO{resultado con personas y el micrófono del prototipo}. \\
\hline
\end{longtable}

% ==============================================================================
\section{Cumplimiento de los objetivos}
% ==============================================================================

\textbf{Objetivo general.} Se desarrolló un sistema de navegación asistida montado en
unos lentes, que combina visión artificial, localización y mapeo simultáneos, trazado de
rayos, aprendizaje por refuerzo e interacción por voz. El sistema funciona de principio a
fin en el prototipo, y el mismo código corre en una computadora portátil con cámara,
con videos grabados y en simulación. \TODO{resumen de cómo se comportó el prototipo en
las pruebas con personas y si se considera cumplido el objetivo general}.

\textbf{Objetivo 1, modelo de visión con precisión mayor al 80\,\%.} Se integraron tres
modelos que detectan y segmentan los objetos relevantes, incluidas las puertas, que el
detector original no reconocía. Las pruebas con video mostraron también sus errores
típicos, como confundir mamparas con puertas o tapetes con hoyos. \TODO{precisión medida
en imágenes etiquetadas del prototipo y veredicto sobre el 80\,\%}.

\textbf{Objetivo 2, reconstrucción del escenario con trazado de rayos.} El mapa del
entorno se construye con Ray Casting sobre una retícula de ocupación, se divide en
bloques, se guarda entre sesiones y se consulta para obtener las distancias que usa la
navegación. Se eligió Ray Casting en lugar de un trazado de rayos completo por su costo,
como se justificó en el marco teórico. El mapa también se corrige cuando el entorno
cambia, porque los ecos precisos liberan las celdas de un obstáculo que se retiró y la
segmentación borra los falsos peligros de piso. El objetivo se cumple.

\textbf{Objetivo 3, navegación con una medición del terreno 90\,\% segura.} En el
simulador, el agente llegó a la meta en el 97.8\,\% de los episodios, chocó en el
2.0\,\% y nunca cayó en un peligro de piso. Con el sistema completo en lazo cerrado,
el usuario virtual llegó en las 20 sesiones y en 17 terminó a menos de un metro del
objeto. Tomando como medida de seguridad la tasa de episodios sin choque ni caída, el
objetivo se cumple en simulación. \TODO{tasa de recorridos seguros con el prototipo y
veredicto sobre el 90\,\%}.

\textbf{Objetivo 4, interacción por voz.} El sistema reconoce comandos de voz sin
conexión, responde en voz alta con una voz natural y se puede usar sin pantalla. El
objetivo se cumple; su punto débil es el desempeño con ruido.

\textbf{Objetivo 5, lenguaje natural para indicaciones cotidianas y de emergencia.}
El PLN reconoce 14 intenciones, entre ellas la emergencia, y la personalidad del
asistente hace que las indicaciones suenen cercanas, con los objetos nombrados en
español. Los comandos de uso frecuente se interpretan bien, pero con frases coloquiales
que no conoce el PLN acierta poco más de la mitad, y algunas emergencias dichas sin
palabras de lesión no se reconocen, así que el objetivo se cumple en parte. \TODO{aviso
a contactos de emergencia}.

% ==============================================================================
\section{Cumplimiento de los requerimientos}
% ==============================================================================

La Tabla~\ref{tab:cumplimiento_req} muestra el estado de cada requerimiento. ``Parcial''
quiere decir que la función existe pero con una limitación medida.

\begin{longtable}{|p{1.6cm}|p{4.6cm}|p{1.9cm}|p{6.4cm}|}
\caption{Estado de los requerimientos funcionales y no funcionales}
\label{tab:cumplimiento_req} \\
\hline
\textbf{Req.} & \textbf{Descripción} & \textbf{Estado} & \textbf{Observación} \\
\hline
\endfirsthead
\hline
\textbf{Req.} & \textbf{Descripción} & \textbf{Estado} & \textbf{Observación} \\
\hline
\endhead
RF-01 & Captura de voz con detección de actividad & Cumplido & Vosk detecta el fin
de cada frase. \\
\hline
RF-02 & Filtrado de audio & No implementado & El filtro de Wiener no mostró mejora y
se dejó fuera. \\
\hline
RF-03 & Reconocimiento e intención & Parcial & 95\,\% de los comandos frecuentes en
limpio, pero 51.2\,\% en frases coloquiales nuevas. \\
\hline
RF-04 & Síntesis de voz & Cumplido & Voz neuronal con respaldo local. \\
\hline
RF-05 & Control de sesión por voz & Cumplido & Iniciar, pausar, reanudar y
detener. \\
\hline
RF-06 & Todas las funciones por voz & Parcial & Las alertas se apagan con tecla o
botón, no por voz. \\
\hline
RF-07 & Captura y preprocesamiento & Parcial & Solo redimensionamiento; los demás
pasos se descartaron. \\
\hline
RF-08 & Detección con localización & Cumplido & Personas, obstáculos, puertas,
escaleras y desniveles. \\
\hline
RF-09 & Clasificación y riesgo & Cumplido & Riesgo de 0 a 3 según la tabla de TT1. \\
\hline
RF-10 & Distancia a los objetos & Cumplido & Monocular, con errores de 25\,\% o más,
y precisa con el ultrasonido dentro de su cono. \\
\hline
RF-11 & Validación con sensores & Cumplido & Validación cruzada con los dos
ultrasónicos frontales. \\
\hline
RF-12 & Alerta de obstáculo & Cumplido & Pitido estéreo y frase corta. \\
\hline
RF-13 & Relaciones contextuales & Parcial & Estado de las puertas, corregido con el
ultrasonido; no se mide la obstrucción de pasillos. \\
\hline
RF-14 & Estimación de la pose & Cumplido & EKF con IMU y respaldo por flujo
óptico. \\
\hline
RF-15 & Mapa incremental & Cumplido & Sin mapa previo. \\
\hline
RF-16 & Zonas libres y ocupadas & Cumplido & Log-odds y capa de peligros. \\
\hline
RF-17 & Actualización dinámica & Cumplido & Los ecos precisos liberan las celdas de
obstáculos retirados. \\
\hline
RF-18 & Ruta local & Cumplido & Agente de PPO y planificador. \\
\hline
RF-19 & Actualización periódica de la ruta & Cumplido & Cada 2~s. \\
\hline
RF-20 & Recálculo ante obstáculos & Cumplido & Decisión inmediata ante una alerta
crítica. \\
\hline
RF-21 & Instrucciones de movimiento & Cumplido & Frases de TT1 y versión
amigable. \\
\hline
RF-22 & Destinos sobre el mapa & Cumplido & Por etiqueta de \emph{landmark}. \\
\hline
RF-23 & Ruta alternativa por voz & Cumplido & Bloquea por tres minutos el tramo
siguiente y replanea; avisa si no hay otra ruta. \\
\hline
RNF-01 & Alerta en menos de 2~s & \TODO{} & Unos 160~ms en la computadora portátil;
\TODO{latencia en la Raspberry Pi}. \\
\hline
RNF-02 & Voz en menos de 500~ms & Cumplido & 10~ms de procesamiento, sin contar la
espera de fin de frase. \\
\hline
RNF-03 & Mapa a 10~Hz & \TODO{} & 1~ms por ciclo para la representación;
\TODO{frecuencia del ciclo completo en la Raspberry Pi}. \\
\hline
RNF-04 & Ruta cada 10~s o menos & Cumplido & Cada 2~s. \\
\hline
RNF-05 & Hardware limitado & \TODO{} & Inferencia de navegación en NumPy;
\TODO{uso de CPU y memoria en la Raspberry Pi}. \\
\hline
RNF-06 & Precisión de detección del 90\,\% & \TODO{} & \TODO{precisión medida}. \\
\hline
RNF-07 & 92\,\% de comandos con 60~dB & No cumplido & 95.0\,\% en limpio y de 62 a
74\,\% con ruido. \\
\hline
RNF-08 & Degradación controlada & Cumplido & Respaldos sin IMU, sin micrófono, sin
voz y sin internet. \\
\hline
RNF-09 & Comportamiento conservador & Cumplido & Capa de seguridad, confianza de la
pose y alertas ``posible'' para detecciones contradichas. \\
\hline
RNF-10 & Interfaz solo por voz & Parcial & Igual que RF-06. \\
\hline
RNF-11 & Instrucciones claras & Cumplido & Frases cortas en español, distancias en
pasos y objetos con su nombre en español. \\
\hline
RNF-12 & Retroalimentación proactiva & Cumplido & Alertas, llegada y cambios de
zona. \\
\hline
RNF-13 & Operación sin internet & Cumplido & Con una voz de respaldo menos
natural. \\
\hline
RNF-14 & Modularidad & Cumplido & Interfaces comunes; el mismo código trabaja con
cámara, video y simulación. \\
\hline
RNF-15 & Actualización de modelos & Cumplido & Modelos intercambiables y script de
reentrenamiento. \\
\hline
\end{longtable}

% ==============================================================================
\section{Conclusiones}
% ==============================================================================

Al terminar TT1, la propuesta era una arquitectura en papel. Al terminar TT2 hay unos
lentes que reciben la imagen de una cámara, entienden una orden hablada, construyen un
mapa del lugar, deciden por dónde ir y lo dicen en voz alta, con todos los módulos
conectados y probados. Ese es el resultado principal del trabajo.

La lección técnica más importante tiene que ver con la forma de evaluar. Varios módulos
parecían terminados hasta que se midieron dentro del sistema completo. El agente de
navegación llega al 97.8\,\% de las metas en su simulador, pero al trabajar junto con el
SLAM y la estimación monocular apareció un problema que solo se nota con todo integrado,
porque la regla de llegada por percepción terminaba las sesiones antes de tiempo. Medir
el sistema en lazo cerrado permitió encontrar la causa, corregirla y pasar de 7 a 17
sesiones de 20 con el usuario a menos de un metro del objeto. Algo parecido pasó con el
EKF, que apenas mejora a la odometría en recorridos cortos y muestra su valor cuando la
odometría tiene sesgo, y con el PLN, que acierta casi siempre las frases parecidas a
las suyas y apenas la mitad de las expresiones coloquiales que nunca vio. Ninguno de estos
hallazgos se habría visto midiendo cada módulo por separado.

La segunda conclusión es que la percepción manda sobre el resto. El mapa, la navegación
y las alertas tardan alrededor de un milisegundo por ciclo, mientras que la visión tarda
más de cien, y con los errores pasa algo similar; una distancia monocular equivocada o un
tapete tomado por hoyo se propagan hasta lo que escucha el usuario. Por eso las mejoras
con más impacto son de percepción, como reentrenar el detector con imágenes del entorno
real o separar la banqueta del arroyo vehicular. En esa misma línea, poner los
ultrasónicos al frente resultó una buena decisión, porque permiten confirmar lo que ve la
cámara, detectar vidrios y mamparas y dar por buena una llegada solo cuando la distancia
está verificada.

La tercera es que vale la pena mantener capas determinísticas por encima de los modelos
aprendidos. La capa de seguridad del navegador, las palabras de seguridad del PLN y el
módulo de alertas siempre encendido no dependen de que una red neuronal acierte. Aun así,
las pruebas mostraron sus límites; en una de las primeras sesiones, un objeto bajo y muy cercano
quedó fuera de la vista de todos los sensores y el sistema siguió indicando avanzar. Para
un sistema que guía a una persona con discapacidad visual, esa redundancia es necesaria,
y también lo es reconocer qué casos no cubre.

Por último, el sistema es un prototipo de investigación, que es el alcance definido para
este trabajo. \TODO{conclusión a partir de las pruebas con el prototipo y de la opinión
de los participantes}. Como se estableció en las limitaciones del proyecto, los lentes
son un complemento del bastón blanco y no lo sustituyen.

% ==============================================================================
\section{Trabajo a futuro}
\label{sec:trabajo_futuro}
% ==============================================================================

El trabajo pendiente se ordena según su efecto sobre la seguridad del usuario y
según lo que se aprendió en las mediciones de esta etapa.

\textbf{Correcciones identificadas en el software.} Varias deficiencias quedaron
localizadas con precisión y su corrección es acotada:

\begin{itemize}
    \item Exigir que la llegada por percepción se cumpla durante varios cuadros
    seguidos, o combinarla con la distancia en el mapa, para que un error monocular
    aislado no termine la sesión.
    \item Agrupar el tiempo de espera de las alertas por tipo y sector, para que una
    precaución lateral no silencie una frontal, y unificar la ortografía de las frases.
    \item Permitir que los rayos precisos reduzcan el log-odds de las celdas
    consolidadas y que la capa de peligros pierda impactos con el tiempo, de modo que
    los obstáculos que desaparecen se borren del mapa (RF-17).
    \item Traducir al español las etiquetas de los \emph{landmarks} en los mensajes de
    llegada y en la descripción del entorno.
    \item Ampliar las frases de referencia del PLN con las expresiones que fallaron en la
    evaluación y revisar la relación entre ``ayuda'' y ``emergencia''.
    \item Hacer que la petición de ruta alternativa obligue al agente a evitar el
    camino anterior.
\end{itemize}

\textbf{Integración del hardware.} Montar los componentes en el armazón diseñado,
conectar la IMU BNO055 y los sensores MB1232 por I2C, medir con regla la geometría
real y ajustarla en \texttt{montaje.py}, ejecutar la calibración de la IMU con los
lentes puestos y calibrar la longitud de paso de cada usuario. Con
los sensores montados se debe medir el ruido real de la odometría, ajustar el
simulador y reentrenar el agente de navegación. También hay que medir los tiempos de
cada módulo en la Raspberry Pi y, si no alcanzan los requerimientos, probar los
modelos de visión en sus versiones más pequeñas o la arquitectura con una computadora
externa que se planteó en TT1.

\textbf{Percepción con datos propios.} Etiquetar imágenes de los pasillos, puertas y
mamparas de la escuela para reentrenar el detector con el script preparado y, con el
mismo conjunto, medir la precisión que exige RNF-06. Separar la banqueta del arroyo
vehicular en la segmentación, para que el sistema no considere transitable la calle.

\textbf{Localización.} Agregar cierre de lazo o una prueba de compatibilidad conjunta
en la asociación de datos para evitar los \emph{landmarks} duplicados y la
divergencia de orientación observada en recorridos largos.

\textbf{Voz en condiciones reales.} Repetir la evaluación de reconocimiento con
personas reales, con el micrófono del prototipo y en lugares con ruido de calle, y
probar un modelo de Vosk más grande si la Raspberry Pi lo permite. Implementar el
aviso a un contacto de emergencia.

\textbf{Pruebas con usuarios.} La validación final del sistema requiere pruebas con
personas con discapacidad visual, en escenarios controlados, acompañadas y con el
bastón blanco en uso. Antes de ello conviene probar el sistema con personas con los
ojos vendados en un entorno conocido, para ajustar el ritmo y la redacción de las
indicaciones.

